\documentclass[11pt]{article}
\usepackage{ochamath}
\subjectcolor{2F5DA8}
\setsheet{線形代数}{最小二乗法　演習}{https://math.ochanote.com/linear-algebra/least-squares/}
\begin{document}
\makesheethead{解答}

\problem[定数モデル]
\begin{ans}
正規方程式は $3c=1+2+6=9$ なので $c=3$。残差は $-2,-1,3$ で和は0、2乗和は $4+1+9=14$。
\end{ans}
\point{定数近似の解は算術平均。}

\problem[直線近似]
\begin{ans}
$A^{\mathsf T}A=\begin{pmatrix}3&3\\3&5\end{pmatrix}$、$A^{\mathsf T}\boldsymbol b=(5,7)^{\mathsf T}$。差を取ると $2c_1=2$ なので $c_1=1,c_0=2/3$。残差は $(1/3,-2/3,1/3)^{\mathsf T}$、2乗和は $2/3$。和と $x$ で重みを付けた和はどちらも0。
\end{ans}
\point{結果は元の測定値との差で確かめる。}
\newpage

\problem[従属する列]
\begin{ans}
出力は $(s,s)^{\mathsf T}$、$s=c_1+c_2$。2乗和は $(1-s)^2+(3-s)^2=2(s-2)^2+2$。最適な出力は $(2,2)^{\mathsf T}$ だが、係数は $c_1+c_2=2$ のすべてで、無数にある。
\end{ans}
\point{射影された出力の一意性と係数の一意性を分ける。}

\problem[QRによる解法]
\begin{ans}
$R=\operatorname{diag}(\sqrt2,1)$、$Q^{\mathsf T}\boldsymbol b=(2\sqrt2,2)^{\mathsf T}$ なので $\boldsymbol c=(2,2)^{\mathsf T}$。出力は $(2,2,2)^{\mathsf T}$、残差は $(-1,0,1)^{\mathsf T}$。$A^{\mathsf T}\boldsymbol r=\boldsymbol0$ で検算。
\end{ans}
\point{長方形QではQの転置でモデル方向の成分を取り出す。}
\end{document}
